\section{总结与延伸}

\subsection{全课知识图谱}

本讲构建了机器人学习的完整认知框架：

\begin{center}
\begin{tikzpicture}[node distance=0.8cm, every node/.style={draw, rounded corners, minimum width=3.5cm, minimum height=0.7cm, font=\small}]
\node (form) {问题建模};
\node (perc) [below=of form] {机器人感知};
\node (rl) [below=of perc] {强化学习};
\node (model) [below left=0.8cm and -0.5cm of rl] {模型学习};
\node (il) [below right=0.8cm and -0.5cm of rl] {模仿学习};
\node (fm) [below=1.6cm of rl] {基础模型};
\draw[->, thick] (form) -- (perc);
\draw[->, thick] (perc) -- (rl);
\draw[->, thick] (rl) -- (model);
\draw[->, thick] (rl) -- (il);
\draw[->, thick] (model) -- (fm);
\draw[->, thick] (il) -- (fm);
\node[draw=none, right=0.8cm of form, text width=5cm, font=\scriptsize] {目标、状态、动作、奖励};
\node[draw=none, right=0.8cm of perc, text width=5cm, font=\scriptsize] {多模态、具身、主动、情境化};
\node[draw=none, right=0.8cm of rl, text width=5cm, font=\scriptsize] {Q-Learning, SAC, PPO};
\node[draw=none, left=0.8cm of model, text width=4cm, font=\scriptsize] {前向模型 + MPC};
\node[draw=none, right=0.8cm of il, text width=4cm, font=\scriptsize] {行为克隆, 扩散策略};
\end{tikzpicture}
\end{center}

\subsection{关键 Takeaways}

\begin{importantbox}{五条核心洞见}
\begin{enumerate}
    \item \textbf{机器人学习 $\neq$ 计算机视觉}：它是序列决策问题，动作影响未来状态，需要处理随机性、延迟奖励和不可微环境
    \item \textbf{运动控制接近解决，操作仍是开放问题}：RL + Sim-to-Real 在 locomotion 中非常成功，但 manipulation 中 Sim-to-Real 差距更大
    \item \textbf{模型学习是``想象力''}：学习环境的前向模型，让机器人能预测动作后果，比纯试错更高效
    \item \textbf{模仿学习是最实用的方法}：通过人类示教快速获得工作策略，扩散策略等新方法大幅提升表达能力
    \item \textbf{基础模型是未来方向}：融合视觉语言知识和大规模机器人数据，追求广泛泛化的通用策略
\end{enumerate}
\end{importantbox}

\subsection{拓展阅读}

\begin{itemize}
    \item Mnih et al., ``Playing Atari with Deep Reinforcement Learning'' (DQN), 2013
    \item Silver et al., ``Mastering the game of Go'' (AlphaGo), Nature 2016
    \item Lee et al., ``Learning quadrupedal locomotion over challenging terrain'', Science Robotics 2020
    \item OpenAI, ``Solving Rubik's Cube with a Robot Hand'', 2019
    \item Chi et al., ``Diffusion Policy'', RSS 2023
    \item Black et al., ``$\pi_0$: A Vision-Language-Action Flow Model for General Robot Control'', 2024
    \item Yunzhu Li AAAI 2025 Tutorial: Robotic Foundation Models
\end{itemize}

\end{document}
